% Compile twice: lualatex 20001210.tex
\documentclass[11pt]{article}
\usepackage[margin=25mm]{geometry}
\usepackage{fontspec,unicode-math}
\setmainfont{Latin Modern Roman}
\setmathfont{Latin Modern Math}
\usepackage{amsmath,amssymb,mathtools}
\usepackage{xurl,hyperref}
\hypersetup{colorlinks=true,urlcolor=blue,linkcolor=blue}
\setcounter{secnumdepth}{0}
\setlength{\parindent}{0pt}
\setlength{\parskip}{5pt}
\setlength{\emergencystretch}{3em}
\newcommand{\sref}[2]{\hyperlink{source:#1:#2}{[S#2]}}
\newcommand{\cataloguescope}[1]{\paragraph{Source question.}#1}
\newcommand{\floor}[1]{\left\lfloor#1\right\rfloor}
\begin{document}
% UnsolvedMath ID 20001210; source catalogue code AIM-COMPUTATION-0048
\setcounter{section}{20001209}
\section{Fourier edge invariants and a visibility threshold for finite rate-class mixtures}\label{20001210}
{\small\sffamily UnsolvedMath ID 20001210 · Computer Science}
\subsection{Definitions and mathematical statement}

\paragraph{Shared notation.}Standard number systems are denoted by $\mathbb N,\mathbb Z,\mathbb Q,\mathbb R,\mathbb C$; index ranges, computational models and specialized notation follow the cited problem source and the definitions below.


A phylogenetic model assigns transition parameters to edges of a tree and maps them to joint probabilities at its observed leaves. A phylogenetic invariant is a polynomial vanishing on all probability vectors in that model. Rate variation among sites can yield mixtures of phylogenetic probability distributions. Algebraically, mixtures are modeled by secant varieties: Zariski closures of spans of a fixed number of points of the original model variety. The task is to find equations for such rate-variable model families.
\paragraph{Statement recorded in the supplied source.}
Problem 34. Find invariants for other models with rate variation among sites: secant varieties of the phylogenetic variaties. \sref{20001210}{2}
\paragraph{Scope qualification.}
The catalogue flags this record as reconstructed and unverified. The supplied question is preserved; the definitions do not certify a unique intended model. The source does not fix the tree, mutation model, number of rate classes, or mixture weights; no one narrower model is substituted for the question.
The supplied catalogue marks this formulation as an unverified reconstruction; the source link does not establish that all missing assumptions were recovered.
\paragraph{Source status.}
The supplied catalogue labels this record partially solved. Its source-status metadata was checked against the supplied snapshot on 2026-10-11. This import does not independently establish current open status or mathematical correctness.
\subsection{Short English statement}
Can polynomial invariants be found for phylogenetic models that allow mutation rates to vary among sites?
\subsection{Sources}
\begin{itemize}
\item[\hypertarget{source:20001210:1}{S1}] Ulam AI and the UnsolvedMath Contributors, UnsolvedMath: A Curated Collection of Open Mathematics Problems, record 20001210 (AIM-COMPUTATION-0048). Catalogue curation and metadata: CC BY 4.0. The supplied source question is retained; definitions and the English summary are editorial additions.
\url{https://huggingface.co/datasets/ulamai/UnsolvedMath}.
\item[\hypertarget{source:20001210:2}{S2}] AIM Workshop Problem Lists, original problem source linked by catalogue record AIM-COMPUTATION-0048
\url{https://aimath.org/WWN/compalgstat/compalgstat.pdf}.
\end{itemize}
\paragraph{Preparation and review.}
Original question attributed to the cited source and UnsolvedMath contributors. Codex (GPT-6; exact serving revision not exposed) assisted with definitions, English summaries, source inspection and formatting. No new human mathematical authorship or independent proof verification is asserted. Canonical ID 20001210 and source code AIM-COMPUTATION-0048 checked against the complete supplied registry and existing repository statements; known flagged repetitions and obvious identical conjectures excluded. Full mathematical-equivalence/literature review has not been performed. Statement preparation and final review followed contribution prompts C and D; automated checks verify structure and identity, not open status.
\label{end:20001210}
\end{document}
